% FILE: 07_open_questions_and_next_work.tex
% Project: research-open-ontologies
% Role: open-ontology-research-and-rdf-authority

\section{Open Questions and Next Work}
\label{sec:research-open-ontologies-open-questions}

\subsection{Known Gaps}
\label{sec:research-open-ontologies-gaps}

Four structural gaps are confirmed by evidence in the research corpus:

\paragraph{Gap 1: TTL Syntax Validation Not Run.}
The conformance audit (PI\_GGEN\_UNIFIED\_RUN\_CONFORMANCE\_AUDIT\_001) Gate~3 FAIL
reports 23 TTL files ``assumed invalid pending tooling installation.'' The \texttt{rapper}
command (\texttt{librdf-utils} package) has never been installed on the research machine.
The Phase~7 roundtrip report confirms that rdflib loads the files without error, but rdflib
is a Python parser and may tolerate syntax that a strict Turtle parser would reject. Until
\texttt{rapper} is run against all 24 TTL files and returns zero errors, this gate cannot
be closed.

\paragraph{Gap 2: SPARQL Prefix Mismatch.}
The Phase~7 roundtrip report shows \texttt{count\_checkpoints=0} and
\texttt{count\_artifacts=0} despite 3{,}684 triples loaded. The \texttt{sample\_subjects}
query confirms that entities \emph{are} present (e.g.,
\texttt{https://process.intelligence/checkpoint/SUBSTRATE\_COMPLETE\_001}). The
\texttt{count\_checkpoints} query must be using a different prefix (likely the class
\texttt{pi:Checkpoint} or \texttt{pi:Artifact} from \path{pi-program.ttl}) than the
instance IRIs actually use. This is a SPARQL namespace alignment gap, not an ontology
design defect. Resolution: inspect the failing queries, align prefixes to the IRI patterns
confirmed by \texttt{sample\_subjects}.

\paragraph{Gap 3: Receipt TTL Absent.}
The Phase~7 report records \texttt{RECEIPT\_NOT\_FOUND} after checking three candidate paths:
\path{receipt.ttl}, \path{pi-ggen-receipt.ttl}, \path{checkpoint-receipt.ttl}. The receipt
chain is JSON-only. The ontology corpus lacks a self-describing provenance graph in RDF
form. This means the ontology layer cannot describe its own provenance using the same
vocabulary it uses to describe other research artifacts, which is an architectural
inconsistency.

\paragraph{Gap 4: Namespace Regime Fragmentation.}
The ontology files use two distinct namespace regimes: \texttt{http://process-intelligence.local/}
(autonomic law, MAPE-K) and \texttt{https://process.intelligence/} (pi-program, compat,
wasm4pm, lifecycle, ma). The declared alignment with public standard vocabularies (PROV-O,
SKOS, ODRL, DCTERMS, SHACL) is declared by prefix import but not verified at the
semantic level. Whether the local class definitions are genuine subclasses or alignments
of the public standard classes is unverified. This is an \textbf{[AUTHOR THESIS / INTERPRETATION]}
gap: the alignment is structurally plausible but has not been validated by a formal ontology
alignment tool.

\subsection{Deferred Gates}
\label{sec:research-open-ontologies-deferred-gates}

The following gates are deferred pending tool availability or resource allocation:

\begin{description}
  \item[Gate 3 (TTL Syntax):]
        Blocked by: \texttt{rapper}/\texttt{librdf-utils} not installed.
        23 violations assumed; actual count unknown.
        Deferred to: next environment provisioning pass.

  \item[Gate 4 (SPARQL Syntax):]
        Blocked by: no \texttt{.rq} query parser installed.
        61 violations assumed; actual count unknown.
        Deferred to: next environment provisioning pass.

  \item[Gate 9 (Legacy .ggen Classified):]
        13 legacy \texttt{.ggen} files in \path{research/pi-program/.ggen/} and
        \path{research/prompt-manufactory/.ggen/} are present but not classified
        in the ggen manifest registry.
        Deferred to: next ggen audit cycle.

  \item[Receipt TTL Manufacture:]
        A receipt TTL must be manufactured that describes the ontology corpus provenance
        in RDF using PROV-O (\texttt{prov:wasDerivedFrom}, \texttt{prov:wasGeneratedBy},
        \texttt{prov:hadPrimarySource}). This artifact would close the architectural
        inconsistency noted in Gap~3.
        Deferred to: next ggen generation rule authoring session.
\end{description}

\subsection{Research Risks}
\label{sec:research-open-ontologies-risks}

\paragraph{Risk 1: rdflib Tolerance Masking Real Errors.}
The 14/14 TTL validation result from rdflib may overstate the health of the corpus. If
\texttt{rapper} finds parse errors in files that rdflib silently accepts, the Phase~7
AVAILABLE verdict would need to be downgraded. This risk is \textbf{medium}: rdflib is
a mature library, but strict Turtle parsers can reject valid-looking files that use
non-standard whitespace or escape sequences.

\paragraph{Risk 2: Namespace Boundary Proliferation.}
The two-regime namespace architecture (\texttt{process-intelligence.local} vs.
\texttt{https://process.intelligence/}) may cause federation failures in future SPARQL
federation scenarios where both regimes are queried together. If the regimes are not
formally aligned as OWL equivalent classes, cross-regime SPARQL queries will silently
return incomplete results.

\paragraph{Risk 3: SPARQL Prefix Mismatch Suppressing Real Data.}
The \texttt{count\_checkpoints=0} and \texttt{count\_artifacts=0} results indicate that
the smoke queries used in \texttt{run-ontology-audit.py} are querying for classes that
do not match the actual IRIs in the corpus. Until this is diagnosed, the audit tool may
report AVAILABLE even when the corpus is missing expected individuals.

\subsection{Next Receipted Work}
\label{sec:research-open-ontologies-next-work}

The following work items are defined as next receipted work for the open-ontologies project,
in priority order:

\begin{enumerate}
  \item \textbf{Install \texttt{rapper} and run syntax validation.} Command:
        \texttt{sudo apt-get install librdf-utils \&\& rapper -i turtle <file>.ttl}
        for each of the 24 TTL files. Expected outcome: zero errors, Gate~3 and Gate~4
        PASS, conformance audit reissued as ALIVE candidate.

  \item \textbf{Fix SPARQL prefix alignment.} Inspect the
        \texttt{count\_checkpoints} and \texttt{count\_artifacts} queries in
        \path{run-ontology-audit.py}. Replace class IRIs with the patterns confirmed by
        \texttt{sample\_subjects} (e.g., \texttt{https://process.intelligence/checkpoint/}).
        Re-run audit; verify non-zero counts.

  \item \textbf{Manufacture receipt.ttl.} Author a new ggen generation rule that
        emits a PROV-O receipt graph describing the ontology corpus provenance:
        which files were loaded, by which operator, at which timestamp, under which
        checkpoint authority. The emitted TTL becomes the fourth SPARQL smoke query
        target and closes the RECEIPT\_NOT\_FOUND gap.

  \item \textbf{Classify legacy .ggen files.} Audit the 13 unclassified legacy
        \texttt{.ggen} files in the pi-program and prompt-manufactory sub-programs.
        For each: determine if it is a standard ggen receipt, a custom extension, or a
        stale artifact. Update the ggen manifest registry accordingly. This closes
        conformance audit Gate~9.

  \item \textbf{Formal ontology alignment verification.} \textbf{[FUTURE WORK]}
        Use an ontology alignment tool (e.g., LogMap, AgreementMakerLight) to verify
        that local class definitions in the corpus are correctly aligned to their
        public standard counterparts in PROV-O, SKOS, and SHACL. This gate is not
        required for program ALIVE but is required for claims of open-standard
        interoperability.
\end{enumerate}
